\subsection{Fiat--Shamir transcript binding for the implementation}
\label{sec:fs-implementation}

We model the Fiat--Shamir execution by replacing each verifier challenge with
the output of a random oracle evaluated on a domain-separated transcript.
For the release parameters,
\[
\mathbb F=\GF(2^{128}),\qquad
\mathbb K=\GF(2^{256}),\qquad
g=16.
\]
Under the current bound \(Q\le 2^{64}\) on random-oracle queries, the target is
\[
\lambda_{\mathrm{FS}}=128+\log_2 Q=192.
\]

The initial transcript absorbs a protocol/version domain separator together
with \(n,k,s\), the public seed, and the checked-setup counter.  At round \(j\),
the prover's current commitment root is absorbed before the fold challenge
\(z_j\in\mathbb K\) is derived.  Sparse commitment levels use the same
sequential transcript: skipped roots are not invented, and skipped fold
challenges are still derived in order from the transcript state fixed by all
previous commitments and challenges.

After the final fold, the verifier checks that the last layer is constant.
The complete last-layer value is absorbed before grinding and before query
derivation.  In the batched protocol, all final component values
\((y_1,\ldots,y_t)\) are absorbed in a fixed canonical order; this is
equivalent to absorbing the full final layer because the verifier separately
checks that every component is constant.

Only then is the grinding seed derived under a grinding-specific domain.
After a valid nonce is found, the nonce is absorbed.  Finally, the \(s\) query
positions are derived under a query-specific domain.

Hence the required transcript order is
\[
\boxed{
\text{parameters/setup}
\rightarrow
\text{roots and fold challenges}
\rightarrow
\text{full last layer}
\rightarrow
\text{grinding}
\rightarrow
\text{query positions}.
}
\]

Conditioned on Merkle binding and on the random-oracle transcript being
collision free, Theorem~6.4 applies round by round exactly as in the
public-coin IOP.  Before \(z_j\) is drawn, \(E_{j-1}\) is fixed by the previous
transcript.  Therefore
\[
\Pr[\text{some bad fold challenge}]
\le
\frac{n(N_1+1/\varepsilon)}{|\mathbb K|}.
\]
Conditioned on all fold challenges being good, the descent leaves an
acceptance witness set of density at most \(1-\delta+n\varepsilon\), and the
query phase contributes
\[
(1-\delta+n\varepsilon)^s.
\]
The implementation therefore targets
\[
\frac{n(N_1+1/\varepsilon)}{|\mathbb K|}
+
(1-\delta+n\varepsilon)^s
\le 2^{-192},
\]
with the \(g\)-bit grinding contribution included in the query term.

\paragraph{Transcript-binding checklist.}
The release implementation must:
\begin{enumerate}
\item use distinct domain separators for setup, fold challenges, grinding,
and query derivation;
\item absorb \(n,k,s\), the seed, and setup counter;
\item absorb every committed root before any dependent challenge;
\item absorb the complete final layer before grinding;
\item absorb the complete final layer before deriving queries;
\item in batching, absorb all \(t\) final values in canonical order;
\item absorb the accepted grinding nonce before deriving queries;
\item derive sparse-schedule skipped fold challenges sequentially from the
same transcript;
\item use unambiguous fixed-width serialization.
\end{enumerate}
